\documentclass[12pt,a4paper]{article}
\usepackage[utf8]{inputenc}
\usepackage[T1]{fontenc}
\usepackage[french]{babel}
\usepackage{geometry}
\usepackage{fancyhdr}
\usepackage{amsmath}
\usepackage{listings}
\usepackage{xcolor}
\usepackage{hyperref}
\usepackage{graphicx}
\usepackage{enumitem}
\usepackage{tcolorbox}

\geometry{margin=2.5cm}
\setlength{\headheight}{14.5pt}

\pagestyle{fancy}
\fancyhf{}
\fancyhead[L]{Correction Examen Blanc 1 - Administration Unix}
\fancyhead[R]{\thepage}
\fancyfoot[C]{}

\lstset{
    language=bash,
    basicstyle=\ttfamily\small,
    keywordstyle=\color{blue},
    commentstyle=\color{gray},
    stringstyle=\color{red},
    numbers=left,
    numberstyle=\tiny\color{gray},
    breaklines=true,
    frame=single,
    showstringspaces=false
}

\definecolor{answerbox}{RGB}{240,255,240}
\newtcolorbox{answerbox}{
    colback=answerbox,
    colframe=green!50!black,
    title=Réponse,
    fonttitle=\bfseries
}

\definecolor{explanationbox}{RGB}{240,248,255}
\newtcolorbox{explanationbox}{
    colback=explanationbox,
    colframe=blue!50!black,
    title=Explication,
    fonttitle=\bfseries
}

\title{\textbf{Correction Examen Blanc 1 - Administration Unix/Linux}\\
\large Corrections et Explications Détaillées}
\author{}
\date{}

\begin{document}

\maketitle
\thispagestyle{fancy}

\section*{Barème}
\begin{itemize}
    \item Question 1 : 10 points
    \item Question 2 : 8 points
    \item Question 3 : 12 points
    \item Question 4 : 10 points
    \item Question 5 : 10 points
    \item Question 6 : 12 points
    \item Question 7 : 10 points
    \item Question 8 : 8 points
    \item \textbf{Total : 80 points}
\end{itemize}

\newpage

% ============================================================
% CORRECTION QUESTION 1
% ============================================================

\section{Correction Question 1 : Hiérarchie des Systèmes de Fichiers (10 points)}

\subsection*{Partie A : Définitions (4 points)}

\begin{answerbox}
\textbf{a) FHS (Filesystem Hierarchy Standard) - 2 points}

Le FHS (Filesystem Hierarchy Standard) est un standard qui définit l'organisation et le contenu des principaux répertoires des systèmes de fichiers Linux/Unix.

\textbf{Rôle :}
\begin{itemize}
    \item Garantit la cohérence entre les différentes distributions Linux
    \item Facilite l'administration système (localisation prévisible des fichiers)
    \item Permet la portabilité des connaissances entre distributions
    \item Définit où placer quels types de fichiers
\end{itemize}

\textbf{Importance :}
\begin{itemize}
    \item Les administrateurs savent où chercher les fichiers de configuration
    \item Les développeurs savent où installer leurs logiciels
    \item Facilite la maintenance et le dépannage
\end{itemize}

\textbf{b) Différence /bin vs /usr/bin - 2 points}

\begin{itemize}
    \item \textbf{/bin} : Exécutables essentiels pour tous les utilisateurs
    \begin{itemize}
        \item Situé dans la partition racine (/)
        \item Disponible même en mode single-user (démarrage minimal)
        \item Contient les commandes de base : ls, cp, mv, bash, etc.
        \item Nécessaire avant le montage de /usr
    \end{itemize}
    \item \textbf{/usr/bin} : Exécutables utilisateur additionnels
    \begin{itemize}
        \item Peut être sur une partition séparée
        \item Monté après la partition racine
        \item Contient les programmes additionnels non essentiels au démarrage
        \item Peut être en lecture seule
    \end{itemize}
\end{itemize}

\textbf{Raison de la distinction :} Cette séparation permet de démarrer le système même si /usr n'est pas monté, et permet de monter /usr en lecture seule pour plus de sécurité.
\end{answerbox}

\subsection*{Partie B : Commandes (3 points)}

\begin{answerbox}
\textbf{a) Lister fichiers avec permissions - 1 point}

\begin{lstlisting}
ls -la /etc
\end{lstlisting}

\textbf{Explication :}
\begin{itemize}
    \item \texttt{-l} : Format long (affiche permissions)
    \item \texttt{-a} : Affiche tous les fichiers (y compris cachés, commençant par .)
\end{itemize}

\textbf{b) Trouver fichiers > 100 Mo - 1 point}

\begin{lstlisting}
find /home -type f -size +100M
\end{lstlisting}

\textbf{Explication :}
\begin{itemize}
    \item \texttt{find /home} : Rechercher dans /home
    \item \texttt{-type f} : Fichiers uniquement (pas répertoires)
    \item \texttt{-size +100M} : Taille supérieure à 100 Mo
\end{itemize}

\textbf{c) Afficher dernières lignes en temps réel - 1 point}

\begin{lstlisting}
tail -f -n 20 /var/log/messages
\end{lstlisting}

ou

\begin{lstlisting}
tail -20f /var/log/messages
\end{lstlisting}

\textbf{Explication :}
\begin{itemize}
    \item \texttt{tail} : Affiche les dernières lignes
    \item \texttt{-n 20} : 20 dernières lignes
    \item \texttt{-f} : Follow (suivre les nouvelles lignes en temps réel)
\end{itemize}
\end{answerbox}

\subsection*{Partie C : Répertoires (3 points)}

\begin{answerbox}
\textbf{/var/log - 1 point}

\begin{itemize}
    \item \textbf{Rôle} : Contient les journaux (logs) du système et des applications
    \item \textbf{Contenu} :
    \begin{itemize}
        \item \texttt{/var/log/messages} : Messages système généraux
        \item \texttt{/var/log/secure} : Logs de sécurité (authentifications)
        \item \texttt{/var/log/httpd/} : Logs du serveur web Apache
        \item \texttt{/var/log/cron} : Logs des tâches planifiées
    \end{itemize}
    \item \textbf{Caractéristique} : Taille variable (peut grandir rapidement)
\end{itemize}

\textbf{/proc - 1 point}

\begin{itemize}
    \item \textbf{Rôle} : Système de fichiers virtuel utilisé comme interface avec le noyau Linux
    \item \textbf{Contenu} :
    \begin{itemize}
        \item \texttt{/proc/cpuinfo} : Informations sur le processeur
        \item \texttt{/proc/meminfo} : Informations sur la mémoire
        \item \texttt{/proc/PID/} : Informations sur un processus spécifique
        \item \texttt{/proc/sys/} : Paramètres du noyau modifiables
    \end{itemize}
    \item \textbf{Caractéristique} : Fichiers virtuels créés dynamiquement en mémoire (pas sur disque)
\end{itemize}

\textbf{/dev - 1 point}

\begin{itemize}
    \item \textbf{Rôle} : Contient les fichiers spéciaux (device files) correspondant aux périphériques
    \item \textbf{Contenu} :
    \begin{itemize}
        \item \texttt{/dev/sda}, \texttt{/dev/sdb} : Disques SATA/SCSI
        \item \texttt{/dev/null} : "Poubelle" (ignore tout ce qui y est écrit)
        \item \texttt{/dev/zero} : Source de zéros
        \item \texttt{/dev/tty} : Terminaux
    \end{itemize}
    \item \textbf{Caractéristique} : Interface entre logiciel et matériel (lecture/écriture = communication avec périphérique)
\end{itemize}
\end{answerbox}

\newpage

% ============================================================
% CORRECTION QUESTION 2
% ============================================================

\section{Correction Question 2 : Types de Fichiers Linux (8 points)}

\subsection*{Partie A : Inodes (4 points)}

\begin{answerbox}
\textbf{a) Définition d'un inode - 2 points}

Un \textbf{inode} (index node) est une structure de données contenant toutes les métadonnées d'un fichier, sauf son nom.

\textbf{Contenu d'un inode :}
\begin{itemize}
    \item \textbf{Type} : Fichier régulier, répertoire, lien symbolique, périphérique, etc.
    \item \textbf{Permissions} : Droits d'accès (rwx) pour propriétaire, groupe, autres
    \item \textbf{Propriétaire} : UID (User ID) et GID (Group ID)
    \item \textbf{Taille} : Taille du fichier en octets
    \item \textbf{Dates} : 
    \begin{itemize}
        \item atime : Dernier accès
        \item mtime : Dernière modification
        \item ctime : Dernier changement de métadonnées
    \end{itemize}
    \item \textbf{Pointeurs vers blocs} : Adresses des blocs de données sur le disque
    \item \textbf{Compteur de références} : Nombre de liens physiques pointant vers cet inode
\end{itemize}

\textbf{b) Nom du fichier - 1 point}

Le nom d'un fichier \textbf{n'est PAS stocké dans l'inode}. Il est stocké dans le \textbf{répertoire parent}, qui associe un nom à un numéro d'inode.

\textbf{Raison :} Cela permet à plusieurs noms (liens physiques) de pointer vers le même inode, et facilite le renommage (changement dans le répertoire parent uniquement).

\textbf{c) Vérifier inodes disponibles - 1 point}

\begin{lstlisting}
df -i
\end{lstlisting}

ou pour un système de fichiers spécifique :

\begin{lstlisting}
df -i /dev/sda1
\end{lstlisting}

\textbf{Explication :} Affiche l'espace inode disponible (IUsed, IFree, IUse\%).
\end{answerbox}

\subsection*{Partie B : Liens (4 points)}

\begin{answerbox}
\textbf{a) Différence hard link vs soft link - 2 points}

\begin{itemize}
    \item \textbf{Lien physique (hard link)} :
    \begin{itemize}
        \item Partage le même inode que le fichier original
        \item Plusieurs noms pointent vers le même inode
        \item Le fichier n'est supprimé que quand tous les liens sont supprimés (compteur = 0)
        \item Doit rester dans le même système de fichiers
        \item Reste valide si le fichier est déplacé (tant que inode existe)
        \item Commande : \texttt{ln cible lien}
    \end{itemize}
    \item \textbf{Lien symbolique (soft link)} :
    \begin{itemize}
        \item A son propre inode séparé
        \item Contient le chemin (absolu ou relatif) vers le fichier cible
        \item Peut traverser les systèmes de fichiers
        \item Devient "cassé" (dangling link) si la cible est supprimée
        \item Commande : \texttt{ln -s cible lien}
    \end{itemize}
\end{itemize}

\textbf{b) Commande lien symbolique - 1 point}

\begin{lstlisting}
ln -s /usr/bin/python3 python
\end{lstlisting}

\textbf{Explication :} Crée un lien symbolique nommé "python" dans le répertoire courant, pointant vers /usr/bin/python3.

\textbf{c) Traversée des systèmes de fichiers - 1 point}

\begin{itemize}
    \item \textbf{Lien symbolique} : \textbf{OUI}, peut pointer vers un fichier sur un autre système de fichiers
    \begin{itemize}
        \item Le lien contient un chemin, pas un inode
        \item Peut donc traverser les montages
    \end{itemize}
    \item \textbf{Lien physique} : \textbf{NON}, doit rester dans le même système de fichiers
    \begin{itemize}
        \item Les inodes sont uniques par système de fichiers
        \item Un inode d'un système de fichiers ne peut pas référencer un inode d'un autre
    \end{itemize}
\end{itemize}
\end{answerbox}

\newpage

% ============================================================
% CORRECTION QUESTION 3
% ============================================================

\section{Correction Question 3 : Démarrage du Système et GRUB (12 points)}

\subsection*{Partie A : Processus de démarrage (6 points)}

\begin{answerbox}
\textbf{Séquence complète du démarrage :}

\begin{enumerate}
    \item \textbf{Power on (Alimentation)} - 0.5 point
    \begin{itemize}
        \item Alimentation du système
        \item Initialisation matérielle de base
    \end{itemize}
    
    \item \textbf{BIOS/UEFI (Firmware)} - 1 point
    \begin{itemize}
        \item POST (Power-On Self-Test) : Vérification matérielle
        \item Détection et initialisation des composants (CPU, RAM, disques)
        \item Recherche du bootloader sur les périphériques de démarrage
    \end{itemize}
    
    \item \textbf{BootLoader (GRUB)} - 1.5 points
    \begin{itemize}
        \item Chargement depuis MBR (BIOS) ou partition EFI (UEFI)
        \item GRUB Stage 1 : Code minimal dans MBR
        \item GRUB Stage 2 : Code complet chargé depuis /boot/grub2/
        \item Affichage du menu de sélection (si configuré)
        \item Chargement du noyau Linux et de l'initramfs en mémoire
        \item Passage du contrôle au noyau avec paramètres de démarrage
    \end{itemize}
    
    \item \textbf{Linux Kernel (Noyau)} - 1.5 points
    \begin{itemize}
        \item Initialisation des structures de données internes
        \item Détection et initialisation du matériel (CPU, mémoire, périphériques)
        \item Chargement des modules nécessaires depuis initramfs
        \item Montage du système de fichiers racine (/)
        \item Libération de l'initramfs de la mémoire
    \end{itemize}
    
    \item \textbf{Processus PID 1 (Init/Systemd)} - 1 point
    \begin{itemize}
        \item Lancement du premier processus utilisateur (PID = 1)
        \item Historiquement : /sbin/init (SysVinit)
        \item Moderne : systemd (remplace init)
        \item Ce processus devient le parent de tous les autres processus
    \end{itemize}
    
    \item \textbf{System Ready (Système prêt)} - 0.5 point
    \begin{itemize}
        \item Démarrage des services système selon leurs dépendances
        \item Activation des targets systemd (multi-user, graphical, etc.)
        \item Système prêt à être utilisé
    \end{itemize}
\end{enumerate}
\end{answerbox}

\subsection*{Partie B : GRUB2 (4 points)}

\begin{answerbox}
\textbf{a) Pourquoi ne pas éditer grub.cfg directement - 1.5 points}

On ne doit \textbf{JAMAIS} éditer directement /boot/grub2/grub.cfg car :

\begin{itemize}
    \item \textbf{Régénération automatique} : Le fichier est généré automatiquement par \texttt{grub2-mkconfig} et toute modification sera écrasée lors de la prochaine régénération
    \item \textbf{Syntaxe complexe} : Risque d'erreur de syntaxe qui pourrait empêcher le démarrage
    \item \textbf{Mises à jour} : Les scripts dans /etc/grub.d/ peuvent être mis à jour, et ces mises à jour seraient perdues
    \item \textbf{Meilleure pratique} : Modifier /etc/default/grub et /etc/grub.d/ puis régénérer
\end{itemize}

\textbf{b) Commande après modification /etc/default/grub - 1 point}

\begin{lstlisting}
sudo grub2-mkconfig -o /boot/grub2/grub.cfg
\end{lstlisting}

\textbf{Explication :}
\begin{itemize}
    \item \texttt{grub2-mkconfig} : Génère la configuration GRUB
    \item \texttt{-o /boot/grub2/grub.cfg} : Spécifie le fichier de sortie
    \item Cette commande lit /etc/default/grub et /etc/grub.d/ pour générer grub.cfg
\end{itemize}

\textbf{c) Initramfs - 1.5 points}

L'\textbf{initramfs} (Initial RAM Filesystem) est un système de fichiers temporaire chargé en mémoire avant le montage de la racine.

\textbf{Rôle :}
\begin{itemize}
    \item Contient les modules du noyau nécessaires pour accéder au disque racine
    \item Contient les outils nécessaires pour monter la racine (mkfs, fsck, etc.)
    \item Exécute des scripts d'initialisation
\end{itemize}

\textbf{Cas nécessaires :}
\begin{itemize}
    \item Disque racine sur \textbf{LVM} (Logical Volume Manager)
    \item Disque racine sur \textbf{RAID}
    \item Système de fichiers \textbf{crypté}
    \item Système de fichiers \textbf{réseau} (NFS)
    \item Tout cas où le noyau a besoin de modules pour accéder à la racine
\end{itemize}

\textbf{Processus :} Initramfs monte la racine, puis passe le contrôle au système réel et peut être libéré de la mémoire.
\end{answerbox}

\subsection*{Partie C : Arrêt du système (2 points)}

\begin{answerbox}
\textbf{a) Arrêt immédiat - 1 point}

\begin{lstlisting}
sudo systemctl poweroff
\end{lstlisting}

ou

\begin{lstlisting}
sudo shutdown -h now
\end{lstlisting}

\textbf{b) Redémarrage programmé - 1 point}

\begin{lstlisting}
sudo shutdown -r +30 "Le système va redémarrer dans 30 minutes pour maintenance"
\end{lstlisting}

\textbf{Explication :}
\begin{itemize}
    \item \texttt{-r} : Redémarrage (reboot)
    \item \texttt{+30} : Dans 30 minutes
    \item Message d'avertissement affiché aux utilisateurs connectés
\end{itemize}
\end{answerbox}

\newpage

% ============================================================
% CORRECTION QUESTION 4
% ============================================================

\section{Correction Question 4 : Gestion des Services (10 points)}

\subsection*{Partie A : Systemd (6 points)}

\begin{answerbox}
\textbf{a) Avantages de systemd - 2 points}

\begin{itemize}
    \item \textbf{Parallélisation} : Démarrage parallèle des services indépendants (plus rapide)
    \item \textbf{Gestion des dépendances} : Dépendances automatiques (requires, wants, after)
    \item \textbf{Configuration déclarative} : Fichiers .service simples et clairs
    \item \textbf{Journal intégré} : journalctl pour tous les logs (format unifié)
    \item \textbf{Activation à la demande} : Services démarrés quand nécessaires (sockets, timers)
    \item \textbf{Meilleure gestion des processus} : Cgroups intégrés
    \item \textbf{Targets} : Remplacement moderne des runlevels
\end{itemize}

\textbf{b) Commandes - 3 points}

\begin{enumerate}
    \item Démarrer httpd :
    \begin{lstlisting}
sudo systemctl start httpd
    \end{lstlisting}
    
    \item Vérifier si actif :
    \begin{lstlisting}
systemctl is-active httpd
    \end{lstlisting}
    Retourne "active" ou "inactive"
    
    \item Activer au démarrage :
    \begin{lstlisting}
sudo systemctl enable httpd
    \end{lstlisting}
    
    \item Logs en temps réel :
    \begin{lstlisting}
journalctl -u httpd -f
    \end{lstlisting}
    \texttt{-u} : unité spécifique, \texttt{-f} : follow (temps réel)
\end{enumerate}

\textbf{c) Différence start vs enable - 1 point}

\begin{itemize}
    \item \texttt{systemctl start} : Démarre le service \textbf{immédiatement} (pour cette session)
    \item \texttt{systemctl enable} : Active le service pour qu'il démarre \textbf{automatiquement au boot} (création de liens symboliques)
    \item \textbf{Important} : On peut faire les deux séparément ou ensemble
    \begin{itemize}
        \item \texttt{start} sans \texttt{enable} : Démarre maintenant mais pas au boot
        \item \texttt{enable} sans \texttt{start} : Configuré pour démarrer au boot mais pas maintenant
    \end{itemize}
\end{itemize}
\end{answerbox}

\subsection*{Partie B : Targets (4 points)}

\begin{answerbox}
\textbf{a) Définition et exemples - 1.5 points}

Un \textbf{target} systemd est un groupe d'unités (services, mounts, etc.) qui définit un état du système. C'est l'équivalent moderne des runlevels.

\textbf{Exemples :}
\begin{itemize}
    \item \texttt{poweroff.target} : Arrêt (runlevel 0)
    \item \texttt{rescue.target} : Mode maintenance (runlevel 1)
    \item \texttt{multi-user.target} : Multi-utilisateurs texte (runlevel 3)
    \item \texttt{graphical.target} : Multi-utilisateurs graphique (runlevel 5)
    \item \texttt{reboot.target} : Redémarrage (runlevel 6)
\end{itemize}

\textbf{b) Passer au mode texte - 1.5 points}

\begin{lstlisting}
sudo systemctl isolate multi-user.target
\end{lstlisting}

\textbf{Explication :} Change le target actif, arrêtant les services du mode graphique et activant ceux du mode texte.

\textbf{c) Target par défaut - 1 point}

\begin{lstlisting}
systemctl get-default
\end{lstlisting}

Retourne le target qui sera activé au démarrage (par défaut : graphical.target ou multi-user.target).
\end{answerbox}

\newpage

% ============================================================
% CORRECTION QUESTION 5
% ============================================================

\section{Correction Question 5 : Gestion des Paquetages (10 points)}

\subsection*{Partie A : Concepts (4 points)}

\begin{answerbox}
\textbf{a) Définition d'un paquetage - 1.5 points}

Un \textbf{paquetage} est une archive contenant :
\begin{itemize}
    \item \textbf{Fichiers} : Binaires, bibliothèques, fichiers de configuration, documentation
    \item \textbf{Métadonnées} : Nom, version, description, dépendances, scripts
    \item \textbf{Scripts} : Pre-install, post-install, pre-remove, post-remove
    \item \textbf{Signature} : Pour vérification d'intégrité et authenticité
\end{itemize}

\textbf{b) Dépendances - 1.5 points}

Les \textbf{dépendances} sont des paquets requis pour qu'un logiciel fonctionne correctement.

\textbf{Exemple :} Un logiciel A peut nécessiter :
\begin{itemize}
    \item Bibliothèque B (dépendance directe)
    \item Bibliothèque C (dépendance directe)
    \item Bibliothèque D requise par C (dépendance transitive)
\end{itemize}

\textbf{Importance :}
\begin{itemize}
    \item \textbf{Résolution automatique} : Le gestionnaire installe automatiquement toutes les dépendances
    \item \textbf{Cohérence} : Garantit que les versions sont compatibles
    \item \textbf{Sécurité} : Les mises à jour de sécurité sont propagées aux dépendances
    \item \textbf{Simplicité} : Une seule commande installe tout
\end{itemize}

\textbf{c) Différence RPM vs YUM - 1 point}

\begin{itemize}
    \item \textbf{RPM} : Outil de bas niveau
    \begin{itemize}
        \item Gère les paquets individuellement
        \item Ne résout PAS les dépendances automatiquement
        \item Nécessite installation manuelle des dépendances
        \item Exemple : \texttt{rpm -i package.rpm}
    \end{itemize}
    \item \textbf{YUM} : Outil de haut niveau
    \begin{itemize}
        \item Surcouche de RPM
        \item Résout automatiquement les dépendances
        \item Télécharge depuis les dépôts configurés
    \end{itemize}
\end{itemize}
\end{answerbox}

\subsection*{Partie B : Commandes YUM (4 points)}

\begin{answerbox}
\textbf{a) Rechercher paquet contenant fichier - 1 point}

\begin{lstlisting}
yum provides /usr/bin/python3
\end{lstlisting}

ou

\begin{lstlisting}
yum provides "*/python3"
\end{lstlisting}

\textbf{b) Installer avec dépendances - 1 point}

\begin{lstlisting}
sudo yum install httpd
\end{lstlisting}

YUM résout et installe automatiquement toutes les dépendances.

\textbf{c) Lister paquets installés - 1 point}

\begin{lstlisting}
yum list installed
\end{lstlisting}

\textbf{d) Vérifier mises à jour - 1 point}

\begin{lstlisting}
yum check-update
\end{lstlisting}

Affiche la liste des paquets ayant des mises à jour disponibles.
\end{answerbox}

\subsection*{Partie C : Dépôts (2 points)}

\begin{answerbox}
\textbf{a) Localisation fichiers dépôts - 1 point}

Les fichiers de configuration des dépôts YUM sont dans :
\begin{lstlisting}
/etc/yum.repos.d/
\end{lstlisting}

Fichiers avec extension \texttt{.repo}

\textbf{b) Lister dépôts actifs - 1 point}

\begin{lstlisting}
yum repolist
\end{lstlisting}

Pour tous les dépôts (actifs et inactifs) :
\begin{lstlisting}
yum repolist all
\end{lstlisting}
\end{answerbox}

\newpage

% ============================================================
% CORRECTION QUESTION 6
% ============================================================

\section{Correction Question 6 : Gestion des Utilisateurs et Groupes (12 points)}

\subsection*{Partie A : Fichiers de configuration (4 points)}

\begin{answerbox}
\textbf{a) Structure /etc/passwd - 1.5 points}

Format : \texttt{login:passwd:uid:gid:comment:home:shell}

\begin{itemize}
    \item \texttt{login} : Nom de connexion
    \item \texttt{passwd} : \textbf{"x"} indique que le mot de passe est dans /etc/shadow
    \item \texttt{uid} : User ID (0 = root, < 1000 = système, >= 1000 = utilisateurs)
    \item \texttt{gid} : Group ID primaire
    \item \texttt{comment} : Nom complet (champ GECOS)
    \item \texttt{home} : Répertoire personnel
    \item \texttt{shell} : Shell par défaut (/bin/bash, /sbin/nologin)
\end{itemize}

\textbf{b) Pourquoi /etc/shadow - 1.5 points}

Les mots de passe sont dans /etc/shadow pour des raisons de \textbf{sécurité} :

\begin{itemize}
    \item \textbf{/etc/passwd} : Permissions 644 (lisible par tous)
    \begin{itemize}
        \item Permettait attaques par dictionnaire
        \item Accessible à tous les utilisateurs
    \end{itemize}
    \item \textbf{/etc/shadow} : Permissions 000 ou 400 (root uniquement)
    \begin{itemize}
        \item Mots de passe cryptés protégés
        \item Contient aussi politiques de validité
        \item Empêche attaques par dictionnaire
    \end{itemize}
\end{itemize}

\textbf{c) Permissions /etc/shadow - 1 point}

\begin{lstlisting}
-r-------- 1 root root /etc/shadow
\end{lstlisting}

\textbf{Permissions :} 400 (lecture root uniquement) ou 000

\textbf{Pourquoi :} Sécurité maximale - seul root peut lire les mots de passe cryptés
\end{answerbox}

\subsection*{Partie B : Création et gestion (5 points)}

\begin{answerbox}
\textbf{a) Créer utilisateur - 2 points}

\begin{lstlisting}
sudo useradd -u 2000 -g users -d /home/etudiant -s /bin/bash -m etudiant
\end{lstlisting}

\textbf{Explication des options :}
\begin{itemize}
    \item \texttt{-u 2000} : UID 2000
    \item \texttt{-g users} : Groupe primaire "users"
    \item \texttt{-d /home/etudiant} : Répertoire personnel
    \item \texttt{-s /bin/bash} : Shell /bin/bash
    \item \texttt{-m} : Créer le répertoire personnel
    \item \texttt{etudiant} : Nom d'utilisateur
\end{itemize}

\textbf{b) Ajouter au groupe wheel - 1.5 points}

\begin{lstlisting}
sudo usermod -aG wheel etudiant
\end{lstlisting}

\textbf{Explication :}
\begin{itemize}
    \item \texttt{-aG} : Ajouter aux groupes secondaires (append to Groups)
    \item \texttt{-a} : Append (sans écraser les groupes existants)
    \item \texttt{-G wheel} : Groupe wheel
\end{itemize}

\textbf{c) Verrouiller compte - 1.5 points}

\begin{lstlisting}
sudo passwd -l utilisateur
\end{lstlisting}

ou

\begin{lstlisting}
sudo usermod -L utilisateur
\end{lstlisting}

\textbf{Explication :} Verrouille le compte (ajoute "!" devant le hash dans /etc/shadow), empêchant la connexion sans supprimer le compte.
\end{answerbox}

\subsection*{Partie C : Permissions (3 points)}

\begin{answerbox}
\textbf{a) Permission 755 - 1 point}

\texttt{755} en mode octal signifie :
\begin{itemize}
    \item \textbf{7 (propriétaire)} : rwx (lecture, écriture, exécution)
    \item \textbf{5 (groupe)} : r-x (lecture, exécution)
    \item \textbf{5 (autres)} : r-x (lecture, exécution)
\end{itemize}

\textbf{Calcul :} r=4, w=2, x=1 → 4+2+1=7, 4+0+1=5

\textbf{b) Différence setuid, setgid, sticky - 1.5 points}

\begin{itemize}
    \item \textbf{Setuid (Set User ID)} :
    \begin{itemize}
        \item Programme exécuté avec privilèges du propriétaire
        \item Exemple : /usr/bin/passwd (modifie /etc/shadow)
        \item Risque de sécurité si mal utilisé
    \end{itemize}
    \item \textbf{Setgid (Set Group ID)} :
    \begin{itemize}
        \item Sur fichiers : Exécution avec privilèges du groupe
        \item Sur répertoires : Nouveaux fichiers héritent du groupe du répertoire
        \item Utile pour collaboration en équipe
    \end{itemize}
    \item \textbf{Sticky Bit} :
    \begin{itemize}
        \item Sur répertoires uniquement
        \item Empêche suppression par non-propriétaires
        \item Exemple : /tmp
    \end{itemize}
\end{itemize}

\textbf{c) Commande permissions rwxr-xr-x - 0.5 point}

\begin{lstlisting}
chmod 755 fichier
\end{lstlisting}

ou

\begin{lstlisting}
chmod u=rwx,g=rx,o=rx fichier
\end{lstlisting}
\end{answerbox}

\newpage

% ============================================================
% CORRECTION QUESTION 7
% ============================================================

\section{Correction Question 7 : Gestion des Disques (10 points)}

\subsection*{Partie A : Partitionnement (4 points)}

\begin{answerbox}
\textbf{a) MBR vs GPT - 2 points}

\begin{itemize}
    \item \textbf{MBR (Master Boot Record)} :
    \begin{itemize}
        \item Ancien standard (1983)
        \item \textbf{Limitations} :
        \begin{itemize}
            \item 4 partitions primaires maximum (ou 3 primaires + 1 étendue)
            \item Taille max : 2.2 To par partition (adressage 32 bits)
            \item Table de partitions dans secteur 0 (512 octets)
        \end{itemize}
        \item Compatible BIOS uniquement
    \end{itemize}
    \item \textbf{GPT (GUID Partition Table)} :
    \begin{itemize}
        \item Standard moderne
        \item \textbf{Avantages} :
        \begin{itemize}
            \item Jusqu'à 128 partitions (standard, extensible)
            \item Taille très grande : jusqu'à 9.4 Zo (adressage 64 bits)
            \item Redondance : Copie de la table
            \item CRC32 : Vérification d'intégrité
        \end{itemize}
        \item Compatible UEFI (natif) et BIOS (avec support)
    \end{itemize}
\end{itemize}

\textbf{b) Commandes - 2 points}

\begin{enumerate}
    \item Lister disques et partitions :
    \begin{lstlisting}
lsblk
    \end{lstlisting}
    
    Avec plus de détails :
    \begin{lstlisting}
lsblk -f
    \end{lstlisting}
    
    \item Obtenir UUID :
    \begin{lstlisting}
sudo blkid
    \end{lstlisting}
    
    ou pour une partition spécifique :
    \begin{lstlisting}
sudo blkid /dev/sda1
    \end{lstlisting}
\end{enumerate}
\end{answerbox}

\subsection*{Partie B : Systèmes de fichiers (3 points)}

\begin{answerbox}
\textbf{a) Journalisation - 2 points}

La \textbf{journalisation} est un mécanisme qui enregistre les modifications avant de les appliquer au système de fichiers.

\textbf{Principe :}
\begin{enumerate}
    \item \textbf{Journal} : Zone spéciale du disque
    \item \textbf{Écriture} : Modifications d'abord écrites dans le journal
    \item \textbf{Commit} : Puis appliquées aux données réelles
    \item \textbf{Récupération} : En cas de crash, rejouer le journal
\end{enumerate}

\textbf{Avantages :}
\begin{itemize}
    \item \textbf{Récupération rapide} : Pas besoin de fsck long après crash
    \item \textbf{Cohérence} : Garantit l'intégrité du système de fichiers
    \item \textbf{Performance} : Moins de vérifications nécessaires
\end{itemize}

\textbf{b) Formater en ext4 - 1 point}

\begin{lstlisting}
sudo mkfs.ext4 /dev/sdb1
\end{lstlisting}

ou

\begin{lstlisting}
sudo mkfs -t ext4 /dev/sdb1
\end{lstlisting}

\textbf{Attention :} Cette opération efface toutes les données de la partition !
\end{answerbox}

\subsection*{Partie C : Montage (3 points)}

\begin{answerbox}
\textbf{a) Rôle /etc/fstab - 1 point}

Le fichier \textbf{/etc/fstab} (File System Table) définit les systèmes de fichiers à monter automatiquement au démarrage.

\textbf{Rôle :}
\begin{itemize}
    \item Configuration statique des montages
    \item Exécuté par \texttt{mount -a} au démarrage
    \item Permet montage automatique sans intervention
\end{itemize}

\textbf{b) Exemple ligne /etc/fstab - 1 point}

\begin{verbatim}
UUID=12345678-1234-1234-1234-123456789abc  /mnt/data  ext4  defaults  0  2
\end{verbatim}

ou avec nom de périphérique :

\begin{verbatim}
/dev/sdb1  /mnt/data  ext4  defaults  0  2
\end{verbatim}

\textbf{Explication des champs :}
\begin{itemize}
    \item \texttt{UUID=...} ou \texttt{/dev/sdb1} : Périphérique
    \item \texttt{/mnt/data} : Point de montage
    \item \texttt{ext4} : Type de système de fichiers
    \item \texttt{defaults} : Options par défaut (rw, suid, dev, exec, auto, nouser, async)
    \item \texttt{0} : dump (sauvegarde) - 0 = pas de sauvegarde
    \item \texttt{2} : fsck (vérification) - 2 = vérifier après la racine
\end{itemize}

\textbf{c) Monter tous les SF de /etc/fstab - 1 point}

\begin{lstlisting}
sudo mount -a
\end{lstlisting}

\textbf{Explication :} Monte tous les systèmes de fichiers définis dans /etc/fstab qui ne sont pas déjà montés.
\end{answerbox}

\newpage

% ============================================================
% CORRECTION QUESTION 8
% ============================================================

\section{Correction Question 8 : Gestion de l'Espace Swap (8 points)}

\subsection*{Partie A : Concepts (4 points)}

\begin{answerbox}
\textbf{a) Définition du swap - 1.5 points}

Le \textbf{swap} est une zone de stockage sur disque utilisée comme extension de la mémoire RAM.

\textbf{Rôle :}
\begin{itemize}
    \item Stocker les pages mémoire inactives quand la RAM est pleine
    \item Permettre d'exécuter plus de processus que la RAM ne le permet
    \item Libérer de la RAM pour les processus actifs
\end{itemize}

\textbf{b) Mécanisme de pagination - 1.5 points}

\begin{enumerate}
    \item \textbf{Pages} : Mémoire divisée en pages (généralement 4 Ko)
    \item \textbf{Pages actives} : En RAM, utilisées fréquemment
    \item \textbf{Pages inactives} : Peu utilisées, candidates pour le swap
    \item \textbf{Algorithme LRU} : Least Recently Used (pages peu utilisées swapées en premier)
    \item \textbf{Page fault} : Accès à une page en swap déclenche son chargement en RAM
    \item \textbf{Thrashing} : Trop de swap = système très lent (pages constamment swapées)
\end{enumerate}

\textbf{c) Partition vs Fichier swap - 1 point}

\begin{itemize}
    \item \textbf{Partition swap} :
    \begin{itemize}
        \item Partition dédiée sur le disque
        \item \textbf{Performance} : Meilleure (accès direct au disque)
        \item \textbf{Taille} : Fixe (nécessite redimensionnement partition pour changer)
    \end{itemize}
    \item \textbf{Fichier swap} :
    \begin{itemize}
        \item Fichier sur un système de fichiers existant
        \item \textbf{Flexibilité} : Taille modifiable facilement (création/suppression)
        \item \textbf{Performance} : Légèrement inférieure (via système de fichiers)
    \end{itemize}
\end{itemize}
\end{answerbox}

\subsection*{Partie B : Configuration (4 points)}

\begin{answerbox}
\textbf{Étapes complètes pour créer un fichier swap de 2 Go :}

\begin{enumerate}
    \item \textbf{Créer le fichier swap (1 point)}
    \begin{lstlisting}
sudo fallocate -l 2G /swapfile
    \end{lstlisting}
    
    Alternative :
    \begin{lstlisting}
sudo dd if=/dev/zero of=/swapfile bs=1M count=2048
    \end{lstlisting}
    (2048 Mo = 2 Go)
    
    \item \textbf{Verrouiller les permissions (0.5 point)}
    \begin{lstlisting}
sudo chmod 600 /swapfile
    \end{lstlisting}
    
    \textbf{Important :} Permissions 600 (lecture/écriture root uniquement) pour sécurité
    
    \item \textbf{Marquer comme swap (0.5 point)}
    \begin{lstlisting}
sudo mkswap /swapfile
    \end{lstlisting}
    
    Crée la signature swap sur le fichier
    
    \item \textbf{Activer le swap (0.5 point)}
    \begin{lstlisting}
sudo swapon /swapfile
    \end{lstlisting}
    
    Active le swap immédiatement
    
    \item \textbf{Rendre permanent (1.5 points)}
    
    Ajouter dans /etc/fstab :
    \begin{verbatim}
/swapfile  none  swap  defaults  0  0
    \end{verbatim}
    
    \textbf{Explication des champs :}
    \begin{itemize}
        \item \texttt{/swapfile} : Fichier swap
        \item \texttt{none} : Pas de point de montage (swap n'est pas monté)
        \item \texttt{swap} : Type swap
        \item \texttt{defaults} : Options par défaut
        \item \texttt{0} : dump (pas de sauvegarde)
        \item \texttt{0} : fsck (pas de vérification)
    \end{itemize}
    
    \textbf{Vérification :}
    \begin{lstlisting}
swapon --show
    \end{lstlisting}
    
    ou
    
    \begin{lstlisting}
free -h
    \end{lstlisting}
\end{enumerate}
\end{answerbox}

\newpage

\section*{Conclusion de la correction}

\textbf{Points importants à retenir :}
\begin{itemize}
    \item Toujours préciser si une commande nécessite \texttt{sudo} ou root
    \item Expliquer les concepts théoriques avec des exemples concrets
    \item Justifier les choix (pourquoi /etc/shadow, pourquoi ne pas éditer grub.cfg, etc.)
    \item Donner des commandes complètes avec toutes les options nécessaires
    \item Expliquer les mécanismes internes (inodes, démarrage, swap, etc.)
\end{itemize}

\textbf{Bon travail !}

\end{document}

